\documentclass[../main.tex]{subfiles}%
\graphicspath{{\subfix{./image/}}{\subfix{../../image/}}}

\begin{document}

\chapter{一般测度论}\label{chapter:ch01}

本章主要是对标准测度论的回顾，尤其关注 $\mathbb{R}^n$ 上的 Radon 测度。

\cref{section:ch01-01}至\cref{section:ch01-04}简述抽象测度论。在\cref{section:ch01-05}中我们建立 Vitali 覆盖定理与 Besicovitch 覆盖定理，后者是\cref{section:ch01-06}与\cref{section:ch01-07}中 Radon 测度的 Lebesgue--Besicovitch 微分定理的关键。\cref{section:ch01-08}给出 Riesz 表示定理的向量值版本。在\cref{section:ch01-09}中我们研究测度序列与函数列的弱紧性。% XREF p=16 "Sections 1.1 through 1.4 are a rapid recounting of abstract measure theory."

读者必要时可查阅\hyperref[chapter:notation]{书末的记号说明}，那里汇总了本书的记号。

\subfile{Chapter-01/section-01.tex}
\subfile{Chapter-01/section-02.tex}
\subfile{Chapter-01/section-03.tex}
\subfile{Chapter-01/section-04.tex}
\subfile{Chapter-01/section-05.tex}
\subfile{Chapter-01/section-06.tex}
\subfile{Chapter-01/section-07.tex}
\subfile{Chapter-01/section-08.tex}
\subfile{Chapter-01/section-09.tex}
\subfile{Chapter-01/section-10.tex}

\end{document}
